% !PRINTED 369
\weberpage{369}
Setzen wir also
\[
\lambda_1 = a_1 t_1 + \frac{b+\sqrt D}{2}\,t_2,
\]
so wird
\[
N(\lambda_1) = a_1 (a_1 t_1^2 + b\,t_1 t_2 + a_2 c\,t_2^2),
\]
und $(a_1,\ b,\ a_2 c)$ ist gleichfalls primitiv. Da nun $\lambda_1$ durch $\mathfrak{a}$ teilbar ist, so muß $a_1$ durch $N(\mathfrak{a})$, d.\,h. durch $a$ teilbar und mithin $=a$ sein. Also ist $\mathfrak{a}$ primär.

Wenn wir in (1) die $\sqrt D$ ein für allemal fest bestimmen, z.\,B. positiv reell oder positiv imaginär, so ist die lineare Form $\lambda$ durch die quadratische Form $T$ \gesperrt{eindeutig bestimmt}. Die Form des konjugierten Ideals würde nicht durch Änderung des Vorzeichens von $\sqrt D$, sondern von $b$ erhalten.

Es seien $\mathfrak{a}$ und $\mathfrak{b}$ irgend zwei zum Führer $Q$ der Ordnung $[Q]$, die wir jetzt mit $O$ bezeichnen wollen, teilerfremde Ideale, und
\weq{5}{
\begin{aligned}
\mu &= \alpha_1 u_1 + \alpha_2 u_2,\\
\nu &= \beta_1 v_1 + \beta_2 v_2
\end{aligned}
}
seien Basisformen von $\mathfrak{a}$ und $\mathfrak{b}$ in $O$. Aus $\mathfrak{a}$ und $\mathfrak{b}$ bilden wir das Produkt
\[
\mathfrak{a}\,\mathfrak{b} = \mathfrak{c}
\]
mit der Basisform in $O$:
\weq{6}{\lambda = \gamma_1 t_1 + \gamma_2 t_2.}
Die zu $\mu$, $\nu$, $\lambda$ gehörigen quadratischen Formen mögen in $U$, $V$, $T$ bezeichnet sein.

1. \gesperrt{Wir nennen $\lambda$ aus $\mu$ und $\nu$, $T$ aus $U$ und $V$ zusammengesetzt oder komponiert.}

Wir setzen:
\weq{7}{
\begin{aligned}
T &= (a,\ b,\ c),\\
U &= (a_1,\ b_1,\ c_1),\\
V &= (a_2,\ b_2,\ c_2),
\end{aligned}
}
und unsere Aufgabe ist gelöst, wenn $T$ aus $U$ und $V$ abgeleitet werden kann.

Ist
\[
\begin{array}{lcl}
\mathfrak{a} & \text{äquivalent mit} & \mathfrak{a}', \\
\mathfrak{b} & \text{\textquotedbl}\quad\text{\textquotedbl} & \mathfrak{b}',
\end{array}
\]
so ist auch
\[
\mathfrak{a}\,\mathfrak{b} \quad \text{\textquotedbl}\quad\text{\textquotedbl} \quad \mathfrak{a}'\mathfrak{b}',
\]
also
\[
\mathfrak{c} \quad \text{\textquotedbl}\quad\text{\textquotedbl} \quad \mathfrak{c}',
\]
und wenn also $U$, $V$ durch äquivalente Formen $U'$, $V'$ ersetzt werden, so tritt auch an Stelle von $T$ eine äquivalente Form $T'$.
